\chapter{Genomic data and experimental framework}
\label{chap:experimental_framework}
% Backward-compatible label for older references.
\label{chap:methods}

Before introducing the learning models themselves, it is useful to
separate the properties of the data from the properties of the
algorithms. This chapter describes the genomic domains used in the
thesis and the rules that determine how information is allowed to flow
through the experimental pipeline.

This separation is especially important in high-dimensional genomic
analysis. Marker filtering, imputation, LD pruning, model selection,
and feature ranking can all introduce optimistic bias if information
from held-out individuals is allowed to influence them. For this
reason, the experimental protocol is treated as part of the method
rather than as a final implementation detail.

\section{Genomic datasets}
\label{sec:genomic_datasets_preprocessing}

The thesis considers two independent livestock genomic domains.

The primary domain is the caprine ADAPTmap dataset. It is used to
develop and evaluate the main experimental pipeline, including breed
classification, sample-efficiency analysis, marker-efficiency
analysis, and comparison of marker-ranking strategies.

A second ovine Sheep HapMap dataset is used for independent
cross-species methodological validation. The ovine data are not
treated as an extension of ADAPTmap and the same SNP identifiers are
not assumed to have the same role across species. Instead, each domain
is processed as a separate classification problem with its own cohort,
label mapping, preprocessing fit, marker panel, and data partition.

\subsection{Primary caprine domain: ADAPTmap}
\label{subsec:adaptmap_dataset}

ADAPTmap was developed to characterize worldwide goat genetic
diversity and population structure
\cite{stella2018adaptmap, colli2018genome}. The source used by the
current implementation is provided in PLINK format and contains
$4653$ sample rows before individual-level eligibility filtering and
$53347$ raw SNP markers.

Breed identity is obtained from the source sample metadata. The
original row identity of each animal is retained throughout the
pipeline so that split membership and preprocessing provenance remain
auditable.

Individual eligibility is determined before the development/test
split. Animals with genotype missingness greater than or equal to
$0.10$ are excluded. Among the remaining individuals, only breeds
represented by at least $30$ eligible animals are retained. These
rules lead to the frozen primary cohort of $2863$ individuals from
$34$ breeds.

The threshold of $30$ individuals is a \emph{cohort-eligibility
rule}. It should not be confused with the sample-size question studied
later in the thesis. RQ2 deliberately constructs a separate controlled
experiment in which the amount of training data per breed is varied.

\subsection{Independent ovine domain: Sheep HapMap}
\label{subsec:sheep_hapmap_dataset}

The Sheep HapMap domain is used to investigate whether the behavior of
the computational framework is specific to ADAPTmap or can be
observed in an independent ovine dataset.

The validation is methodological rather than marker-level. A marker
ranked highly in goats is not expected, by construction, to
correspond to the same biological locus in sheep. What can be compared
across the two domains are broader properties: whether compact
representations support breed classification, how performance changes
under data scarcity, and whether reduced panels retain useful
discriminative information.

The ovine domain is the ISGC SNP50 HapMap and Sheep Breed Diversity
dataset \cite{kijas2012sheep}, distributed by the CSIRO Data Access
Portal \cite{csiroSheepHapMap2013}. The audited Breedv1 release used
here contains $2819$ animals from $74$ breeds genotyped on the ovine
SNP50 BeadChip, with $49{,}034$ markers retained after the consortium
quality filtering. The distributed PED/MAP files were converted to
PLINK binary with \texttt{plink --sheep --make-bed}.

Individual eligibility follows the same pre-specified eligibility rule as the primary
domain: animals with genotype missingness greater than or equal to
$0.10$ are excluded. All $2819$ animals pass this filter (maximum
missingness $1.99\%$). Among the remaining individuals, breeds
represented by at least $30$ eligible animals are retained, yielding a
frozen ovine cohort of $28$ breeds and $1863$ animals. The cohort is
partitioned into $1585$ development and $278$ held-out test individuals
using the same $85/15$ fractions and the same seeds ($42/42$) as the
caprine domain.

\begin{figure}[ht]
  \centering
  \includegraphics[width=\textwidth]{fig_dataset_cohort}
  \caption{Population structure and cohort composition. For each domain,
  the left panel shows a PCA of the fold-0 training dosages (LD-pruned
  panel; each color is one breed) and the right panel the distribution
  of cohort sizes across breeds together with the $\geq 30$ eligibility
  threshold.}
  \label{fig:dataset_cohort}
\end{figure}

\section{From raw genotypes to model input}
\label{sec:genotype_preprocessing}

Raw genotype matrices require quality control before they can be used
by the learning models. The purpose of preprocessing in this thesis is
not to optimize a score using all available individuals, but to
estimate every data-dependent transformation only from the training
portion that is legitimately available at that stage.

Common genomic quality-control operations include sample missingness,
marker missingness, Minor Allele Frequency (MAF) filtering, and LD
pruning \cite{purcell2007plink, anderson2010data,
chang2015second}. The frozen pipeline applies the following
marker-level sequence:

\begin{enumerate}

    \item markers with training-set missingness greater than or equal
    to $0.10$ are removed;

    \item MAF is computed using the observed training genotypes only,
    and markers with $\mathrm{MAF}<0.01$ are removed;

    \item missing values are imputed using the genotype mode estimated
    from the training data for each retained marker; if multiple
    genotype states are tied, the lowest state is selected
    deterministically;

    \item greedy LD pruning is performed in source marker order using a
    window of $50$ markers, removing a marker when $r^2>0.2$ with a
    marker already retained in the active window.

\end{enumerate}

No Hardy--Weinberg-equilibrium filter is applied and the allele
dosages are not standardized.

Once fitted, the preprocessing state is applied unchanged to held-out
individuals. The post-QC representation, therefore, contains integer
allele dosages

\begin{equation}
    g_{i,j}\in\{0,1,2\},
    \label{eq:allele_dosage}
\end{equation}

with no missing values. This matrix is the common input from which the
variational and contrastive models construct their model-specific
representations.

\section{Development data and the held-out test set}
\label{sec:data_partitioning}

The $2863$ eligible ADAPTmap individuals are divided into two
different roles.

The \emph{development set} contains $2437$ individuals. All model
configuration, cross-validation, preprocessing decisions,
training-control choices, and marker-ranking development must take
place inside this set.

The remaining $426$ individuals form the \emph{held-out test set}. It is
excluded from model development and is opened only once, after the
configuration of the final models has already been fixed. This
distinction is important: a test set can provide an unbiased estimate
of final generalization only if its outcomes are not used to make
subsequent modeling decisions \cite{ambroise2002selection}.

For development, three fixed stratified folds are used:

\begin{center}
\begin{tabular}{c|cc}
\hline
Fold & Training & Validation \\
\hline
0 & 1624 & 813 \\
1 & 1625 & 812 \\
2 & 1625 & 812 \\
\hline
\end{tabular}
\end{center}

For each fold, preprocessing is fitted again using only that fold's
training individuals. As a consequence, the retained SNP panel is
allowed to differ slightly from one fold to another. This is
preferable to building a global marker intersection or union, because
such a global construction would mix information across training
partitions.

In the frozen primary protocol, the three fold-specific post-QC
panels contain $46442$, $46402$, and $46444$ SNPs. These values are
properties of the fold-specific preprocessing fit, not fixed
dimensions imposed on the raw data.

\section{A common label space}
\label{sec:canonical_labels}

For a task containing $C$ breeds, textual breed identifiers are
mapped deterministically to dense class indices

\begin{equation}
    y_i\in\{0,\ldots,C-1\}.
\end{equation}

The same canonical mapping is used by all models evaluated on the same
task cohort. For the primary ADAPTmap classification task, the frozen
cohort contains $34$ breeds. Throughout the thesis, breed names are used in
the narrative text, while the original ADAPTmap identifiers are
retained when compact labels are required, for example, in confusion
matrices or computational artifacts. The correspondence is reported
in Table~\ref{tab:adaptmap_breed_mapping}.

\begin{table}[ht]
\centering
\caption{Breed names and ADAPTmap identifiers for the primary
34-breed caprine cohort.}
\label{tab:adaptmap_breed_mapping}
\begin{tabular}{ll@{\qquad}ll}
\hline
Breed & ID & Breed & ID \\
\hline
Abergelle & ABR &
Alpine & ALP \\

Angora & ANG &
Boer & BOE \\

Bari & BRI &
Barki & BRK \\

Burundi goat & BUR &
Bugituri & BUT \\

Cameroon goat & CAM &
Canindé & CAN \\

Cashmere & CAS &
Creole & CRE \\

Corse & CRS &
Girgentana & GGT \\

Gumez & GUM &
Kamori & KAM \\

Keffa & KEF &
Kilis & KLS \\

Landim & LND &
Landrace & LNR \\

Malagueña & MLG &
Moxotó & MOX \\

Nubian & NBN &
Oasis & OSS \\

Pateri & PAT &
Rangeland & RAN \\

Rossa Mediterranea & RME &
Saanen & SAA \\

Sarda & SAR &
Small East African & SEA \\

Saidi & SID &
Teddi & TED \\

West African Dwarf & WAD &
Woyito Guji & WYG \\
\hline
\end{tabular}
\end{table}

Using a common label order is more than an implementation convenience.
It makes fold-level per-class recall and confusion matrices directly
comparable and allows them to be interpreted without remapping class
indices after every experiment.

\section{The development questions as controlled experiments}
\label{sec:experimental_design}

The thesis studies several questions using the same general data
framework. The difference between experiments lies in what is
deliberately changed while the remaining protocol is held fixed.

For breed classification, the complete post-QC genomic representation
is used as the reference setting. For sample efficiency, the selected
VMGP configuration is held fixed while the number of training
individuals per breed is reduced; the corresponding 15-breed
validation population remains unchanged within each cross-validation fold.
For marker efficiency, the input panel is progressively reduced using
a ranking constructed from the training partition only.
Finally, different marker-ranking methods are compared by evaluating the downstream panels that they
produce.

This view is useful because it avoids treating every experiment as an
independent benchmark. They are different probes of the same
underlying question: how much information is required to preserve
breed-discriminative structure?

\begin{figure}[htbp]
\centering
\begin{tikzpicture}[
    >=Latex,
    font=\small,
    box/.style={draw, semithick, rounded corners=2pt, align=center,
                minimum height=10mm, minimum width=39mm, inner sep=4pt, fill=black!2},
    key/.style={box, fill=black!7},
    small/.style={draw, semithick, rounded corners=2pt, align=center,
                  minimum height=9mm, minimum width=34mm, inner sep=3pt, fill=black!2},
    note/.style={draw, dashed, semithick, rounded corners=2pt, align=center,
                 minimum height=9mm, minimum width=27mm, inner sep=3pt},
    arr/.style={->, semithick},
    darr/.style={->, semithick, dashed}
]
\node[key] (raw) at (0,5.0) {Raw genotype data};
\node[box] (qc) at (0,3.65) {Quality control + encoding};
\node[key] (split) at (0,2.3) {Train / validation / test split};
\node[box] (train) at (-2.9,0.85) {Training partition};
\node[box] (test) at (3.2,0.85) {Held-out test partition};

\node[note] (sub) at (-6.15,0.85) {Subsample $n$\\per breed};
\node[small] (repr) at (-3.8,-0.85) {Representation learning};
\node[small] (classic) at (0.0,-0.85) {Classical ranking};
\node[small] (xai) at (-3.8,-2.25) {XAI attribution};
\node[key] (rank) at (-1.9,-3.65) {Training-only SNP ranking};
\node[key] (panel) at (-1.9,-5.05) {Reduced marker panels};
\node[box] (clf) at (-1.9,-6.45) {Downstream classifier};
\node[key] (eval) at (3.2,-6.45) {Final test evaluation};

\draw[arr] (raw)--(qc); \draw[arr] (qc)--(split);
\draw[arr] (split)--(train); \draw[arr] (split)--(test);
\draw[arr] (train)--(repr); \draw[arr] (train)--(classic);
\draw[darr] (train.west) -- (sub.east);
\draw[arr] (repr)--(xai); \draw[arr] (xai)--(rank); \draw[arr] (classic)--(rank);
\draw[arr] (rank)--(panel); \draw[arr] (panel)--(clf); \draw[arr] (clf)--(eval);
\draw[arr] (test) -- (eval);
\end{tikzpicture}

\caption{Leakage-safe experimental design. Model fitting and SNP ranking are performed using training data, whereas the held-out test partition is reserved for final evaluation. The dashed branch represents the repeated per-breed subsampling used in sample-size experiments. Author's elaboration.}
\label{fig:experimental_design}
\end{figure}


\section{Cross-validation and model selection}
\label{sec:rq1_protocol}

For RQ1, scalar classification metrics are computed on each of the
three cross-validation folds. The primary model-selection metric is
macro-averaged F1.

Configurations are selected using the following fixed rule:

\begin{enumerate}
    \item maximize mean Macro-F1 across the three folds;
    \item if tied, maximize mean Balanced Accuracy;
    \item if still tied, maximize mean Accuracy;
    \item if an exact tie remains, preserve deterministic
    configuration order.
\end{enumerate}

Validation loss is used only as a training-control signal for early
stopping and checkpoint restoration. It is not used as the scientific
ranking metric between candidate model configurations.

The important consequence is that there is no ``best fold'' that is
reported as the final cross-validation result. A configuration is
characterized by its behavior across all three held-out development
partitions.

\section{Studying sample scarcity}
\label{sec:sample_size_protocol}

RQ2 asks how breed-classification performance changes when fewer
training individuals are available for each breed. Unlike RQ1, the
objective is no longer to compare alternative representation-learning
approaches. The sample-efficiency experiment, therefore, uses the
VMGP-inspired model selected during RQ1 as a fixed reference model,
with latent dimensionality $d=96$ and all architectural, loss, and
optimization choices held fixed.

This distinction is important because the experiment is intended to
isolate the effect of training-population size. Repeating model
selection independently at every sample size would introduce an
additional source of variation and would make it difficult to
determine whether changes in performance were caused by the amount of
available data or by changes in model configuration.

A controlled comparison also requires the classification task itself
to remain unchanged as the training population is reduced. The RQ2
cohort is therefore restricted to breeds for which at least $30$
training individuals are available in every pre-specified development
fold. The value $30$ corresponds to the largest per-breed training
population considered in the experiment.

Applying this requirement yields a fixed cohort of 15 breeds:
Abergelle, Alpine, Angora, Boer, Barki, Burundi goat, Creole,
Landrace goat, Nubian, Oasis, Rangeland, Saanen, Small East African,
Saidi, and West African Dwarf. Their corresponding ADAPTmap
identifiers are reported in
Table~\ref{tab:adaptmap_breed_mapping}.

Using the same breed set for every sample-size condition is essential
for interpretation. If breeds entered or disappeared from the
classification problem as the number of available individuals changed,
differences in performance could reflect changes in class composition
or task difficulty rather than changes in training-population size.

The evaluated numbers of training individuals per breed are

\begin{equation}
    \mathcal{N}
    =
    \{5,10,15,20,25,30\}.
    \label{eq:sample_sizes}
\end{equation}

For a given fold and a given $N\in\mathcal{N}$, exactly $N$ training
individuals are sampled from every breed. The resulting training
population therefore contains $15N$ animals, ranging from $75$
individuals for $N=5$ to $450$ individuals for $N=30$.

The validation population is not subsampled. For each development
fold, all validation individuals belonging to the same 15 breeds are
retained, and their source identities remain unchanged across all
sample sizes and sampling replicates. Consequently, differences along
the learning curve are evaluated against the same validation
population within a fold.

A new genomic preprocessor is fitted independently for every
combination of cross-validation fold, sampling replicate, and sample size.
Only the selected $15N$ training individuals are used to estimate its
state. The fitted preprocessor is then applied to both the sampled
training population and the unchanged validation population.

This choice is deliberate because sample scarcity affects more than
neural-network optimization. Training-dependent genomic quantities,
including marker missingness, minor-allele frequency, genotype-mode
imputation, and linkage-disequilibrium pruning, are themselves
estimated from the available animals. Reusing a preprocessing state
estimated from the complete fold would therefore give the smaller
sample-size conditions information derived from animals that are
supposed to be unavailable.

RQ2 consequently measures \emph{end-to-end sample scarcity}: reducing
$N$ changes both the information available for genomic preprocessing
and the number of observations available for representation learning.

The particular animals selected at a given value of $N$ provide
another potential source of variability. Individuals belonging to the
same breed are not necessarily interchangeable: within-breed genetic
diversity, relatedness, population substructure, and less common
genotype profiles can make some subsets more informative than others.
A learning curve obtained from only one random subsample could
therefore partly reflect an accidental choice of animals rather than
the general effect of sample size.

To reduce this dependence, the sampling procedure is repeated using
five pre-specified random seeds,

\begin{equation}
    \mathcal{R}
    =
    \{42,43,44,45,46\}.
    \label{eq:rq2_replicate_seeds}
\end{equation}

For every cross-validation fold and replicate, a deterministic random
ordering of the available training individuals is generated separately
for each breed. Sample-size conditions are derived from the same
ordering in nested form,

\begin{equation}
S_5
\subset
S_{10}
\subset
S_{15}
\subset
S_{20}
\subset
S_{25}
\subset
S_{30}.
\label{eq:rq2_nested_samples}
\end{equation}

Thus, moving from $N=5$ to $N=10$ within a replicate adds five
individuals per breed to the existing population rather than replacing
the original five with an unrelated sample. The same principle applies
to all subsequent values of $N$.

Five sampling replicates are used as a compromise between
characterizing the variability induced by training-set composition and
the computational cost of repeatedly fitting the complete genomic
pipeline. The number five is not intended as a universal statistical
requirement. The purpose of repeated subsampling is instead to avoid
drawing conclusions from a single arbitrary selection of animals.

Together with the three pre-specified cross-validation folds, the five
sampling replicates provide

\begin{equation}
    3 \times 5 = 15
\end{equation}

performance measurements for each value of $N$. These measurements
are used descriptively to characterize average performance and its
variability. They should not be interpreted as 15 fully independent
experimental populations, because measurements within the same
development framework may share validation individuals and partially
overlapping training subsets.

One training detail differs intentionally from the primary RQ1
protocol. RQ2 uses mini-batches of at most $64$ individuals but retains
the final incomplete mini-batch rather than discarding it. This ensures
that every individual selected for a sample-size condition can
contribute to optimization. Discarding the final mini-batch would
remove a different proportion of the available training population for
different values of $N$, which would be undesirable when the number of
available individuals is itself the experimental variable.

Throughout this experiment, the held-out test set remains completely
outside the RQ2 data path. It is neither used for sampling nor
transformed by any RQ2-specific preprocessing state.

\section{Evaluation metrics}
\label{sec:evaluation_metrics}

No single metric describes every aspect of a multiclass classifier.
The primary metric used in the thesis is macro-averaged F1,

\begin{equation}
    F_{1}^{\mathrm{macro}}
    =
    \frac{1}{C}
    \sum_{c=1}^{C}
    F_{1,c},
    \label{eq:macro_f1}
\end{equation}

which gives equal weight to each breed \cite{sokolova2009systematic}.

Balanced Accuracy is reported as the mean recall across classes,

\begin{equation}
    \operatorname{BAcc}
    =
    \frac{1}{C}
    \sum_{c=1}^{C}
    \operatorname{Recall}_c.
\end{equation}

Standard multiclass Accuracy is included as a secondary measure.

For cross-validation, scalar metrics are reported using the mean and
standard deviation across the three frozen folds. Per-class recall and
the raw confusion matrix are retained as diagnostic information.
Because the canonical breed mapping is identical across folds,
class-specific quantities remain directly interpretable.

The RQ2 sample-efficiency analysis additionally includes the five
pre-specified subsampling replicates. For each value of
$N\in\mathcal{N}$, the three cross-validation folds and five sampling
replicates therefore yield 15 performance measurements.

Macro-F1, Balanced Accuracy, and Accuracy are summarized using their
arithmetic mean and sample standard deviation across these 15
measurements. Macro-F1 remains the primary quantity used to interpret
the sample-efficiency curve and to determine the operational minimum
training-population size.

\section{Defining ``minimum'' before looking at the curves}
\label{sec:minimum_requirements}

The terms ``minimum number of individuals'' and ``minimum number of
markers'' can become ambiguous if the threshold is chosen only after
looking at experimental curves. The thesis, therefore, uses an
operational retained-performance definition.

Let $S_{\mathrm{full}}$ be the Macro-F1 obtained using the complete
post-QC marker panel, and let $S(P')$ be the score obtained with a
reduced panel containing $P'$ markers. A minimum retained panel can be
defined as

\begin{equation}
    P_{\min}
    =
    \min
    \left\{
        P':
        S(P')
        \geq
        (1-\epsilon_P)
        S_{\mathrm{full}}
    \right\}.
    \label{eq:minimum_marker_definition}
\end{equation}

Similarly, if $S(N')$ denotes the score obtained using $N'$ training
individuals per breed and $S_{\mathrm{reference}}$ is the pre-specified
reference score,

\begin{equation}
    N_{\min}
    =
    \min
    \left\{
        N':
        S(N')
        \geq
        (1-\epsilon_N)
        S_{\mathrm{reference}}
    \right\}.
    \label{eq:minimum_sample_definition}
\end{equation}

For RQ2, the reference performance is defined internally within the
fixed 15-breed sample-efficiency task. Let

\begin{equation}
    S_{\mathrm{ref}}
    =
    \overline{F_1^{\mathrm{macro}}}(30)
\end{equation}

denote the mean Macro-F1 obtained at $N=30$ across the three
cross-validation folds and five subsampling replicates.

Before inspecting the definitive sample-efficiency curve, the
retained-performance tolerance was fixed as

\begin{equation}
    \epsilon_N = 0.02.
\end{equation}

The operational definition used in RQ2 is consequently

\begin{equation}
    N_{\min}
    =
    \min
    \left\{
        N\in\mathcal{N}:
        \overline{F_1^{\mathrm{macro}}}(N)
        \geq
        0.98\,
        \overline{F_1^{\mathrm{macro}}}(30)
    \right\}.
    \label{eq:rq2_minimum_sample_definition}
\end{equation}

The $98\%$ criterion refers to retained performance relative to the
$N=30$ reference condition; it does not require an absolute Macro-F1
of $0.98$. Fixing this rule before examining the definitive RQ2 curve
prevents the definition of ``minimum'' from being adjusted
retrospectively to match the observed results.

The tolerances must be fixed before the definitive curves are
interpreted. In the frozen protocol they are

\begin{equation}
    \epsilon_N = 0.02,
    \qquad
    \epsilon_P = 0.02 \;\text{(primary)},
    \qquad
    \epsilon_P = 0.05 \;\text{(secondary)}.
    \label{eq:frozen_tolerances}
\end{equation}

The marker-efficiency grid is
$K\in\{50,100,200,500,1000,2000,5000\}$. For the ovine
sample-efficiency mini-sweep, where the minimum fold-training count per
breed is $17$, the grid is truncated to $N\in\{5,10,15\}$ with $N=15$
as the internal reference; the ovine $N_{\min}$ is therefore bounded
above by the reference value and is interpreted through the shape of
the retention curve, not as an extended minimum.
